% !TeX program = pdflatex
%==============================================================================
%  Lista de Revisão --- Aulas 01 a 06 (Bloco I)
%  Material de preparação para a Avaliação Teórica 01.
%
%  O gabarito sai de "Lista de revisao 01-06 - gabarito.tex", que apenas liga a
%  opção [gabarito] e inclui este arquivo --- fonte única, dois PDFs.
%
%  O CORPO DE CADA EXERCÍCIO ESTÁ EM exercicios/ex-NN-slug.tex, e é incluído por
%  \input também pela "Avaliacao teorica 01.tex". Assim é impossível que a prova
%  e a lista divirjam num exercício sorteado: os dois leem o mesmo arquivo.
%  A pontuação por item, que só a prova mostra, vem do comando \pts do estilo.
%
%  Os treze exercícios são inéditos: nenhum repete os vinte das
%  Lista teorica 01--06. Vários pedem derivações que as notas enunciam sem
%  demonstrar (o atalho do LOOCV, as fórmulas do caso ortonormal, as duas formas
%  da fórmula de k dobras).
%==============================================================================
\providecommand{\gabaritoopt}{}
\documentclass[11pt,a4paper]{article}
\usepackage[\gabaritoopt]{../recursos/latex/estilo-avaliacao}

\begin{document}
\cabecalhorevisao{01 a 06}{Fundamentos, Regressão, Validação Cruzada, KNN,
Árvores e \emph{Pipelines}}

\noindent
Treze exercícios para revisar o Bloco~I antes da avaliação. Eles não repetem os
das listas de cada aula --- a maior parte pede que você \emph{demonstre}
resultados que as notas usaram sem demonstrar, ou que ligue duas aulas vistas
separadamente.

\medskip\noindent
Não há ordem obrigatória, mas os exercícios de uma mesma aula se apoiam uns nos
outros, e vários apontam para exercícios das listas anteriores --- quando isso
acontece, o enunciado avisa.

\medskip\noindent
Salvo indicação em contrário, $(\X,Y)$ é um par aleatório com distribuição $P$,
$r(\x)=\E[Y\mid\X=\x]$ é a função de regressão, $\Dados=\{(\X_i,Y_i)\}_{i=1}^n$ é
a amostra de treino (i.i.d.), $p$ é o número de covariáveis e a perda é a
quadrática.

%==============================================================================
\section*{Aula 01 --- Introdução}
%==============================================================================

\begin{exercicio}[o melhor preditor de todos, e o que sobra depois dele]
\label{ex:01-ropt}
% Aula 01 --- o preditor ótimo e o erro irredutível.
% Incluído por "Lista de revisao 01-06.tex" e, se sorteado, pela prova.
No início do curso, vimos uma afirmação forte: de todas as funções que você poderia usar
para prever $Y$ a partir de $\X$, a melhor é a função de regressão
$r(\x)=\E[Y\mid\X=\x]$. Não a melhor entre as lineares, nem a melhor entre as que
o seu método consegue produzir --- a melhor, ponto.

Uma afirmação dessas pede demonstração, e ela é mais curta do que parece. Todo o
trabalho está num truque que você vai reencontrar várias vezes no curso: somar e
subtrair a quantidade certa dentro de um quadrado.
\begin{itens}
  \item \pts{1,2} Seja $g:\R^p\to\R$ uma função qualquer com
        $\E[g(\X)^2]<\infty$. Some e subtraia $r(\X)$ dentro do quadrado e mostre
        que
        \[
          \E\big[(Y-g(\X))^2\big]
            = \E\big[(Y-r(\X))^2\big] + \E\big[(r(\X)-g(\X))^2\big].
        \]
        O trabalho está em mostrar que o termo cruzado desaparece; condicione em
        $\X$ e pergunte-se o que $\E[Y-r(\X)\mid\X]$ vale.
  \item \pts{0,9} Olhe para a identidade que você acabou de provar e responda: por
        que ela mostra que $r$ é o minimizador? Em seguida, mostre que o valor
        mínimo é $\risco(r)=\E[\Var(Y\mid\X)]$.
  \item \pts{0,8} Esse mínimo tem nome: \textbf{erro irredutível}. O nome sugere
        que não há como baixá-lo, e é verdade --- mas a razão não é que os métodos
        atuais sejam ruins. Explique de onde vem essa impossibilidade, e diga o que
        teria de mudar \emph{no problema} para o número diminuir.
  \item \pts{1,1} Na Lista Teórica~01, mostramos que o melhor preditor
        \emph{constante} tem risco $\Var(Y)$. Use a identidade de (a) para escrever
        a diferença $\Var(Y)-\risco(r)$ como uma esperança, e interprete: que
        pergunta prática essa quantidade responde?
\end{itens}

\begin{solucao}
Suponha, ao longo de todo o exercício, que $\E[Y^2]<\infty$: é o que garante que
todas as esperanças abaixo existam. Em particular $r(\X)$ também tem segundo momento
finito, porque $\E[r(\X)^2]=\E\big[(\E[Y\mid\X])^2\big]\le\E\big[\E[Y^2\mid\X]\big]=\E[Y^2]$
(desigualdade de Jensen, aplicada condicionalmente a $\X$).

\smallskip
\textbf{(a)} Some e subtraia $r(\X)$ e expanda o quadrado:
\[
  (Y-g(\X))^2 = \big(Y-r(\X)\big)^2
     + 2\big(Y-r(\X)\big)\big(r(\X)-g(\X)\big)
     + \big(r(\X)-g(\X)\big)^2 .
\]
As três parcelas são integráveis: as duas quadráticas porque $Y$, $r(\X)$ e $g(\X)$
têm segundo momento finito, e o termo cruzado pela desigualdade de Cauchy--Schwarz,
$\E\abs{UV}\le\sqrt{\E[U^2]\,\E[V^2]}$. Podemos, então, tomar esperança termo a
termo, e tudo se resume a mostrar que o termo cruzado tem esperança zero.

Condicione em $\X$ e use a lei da esperança total:
\[
\begin{aligned}
  \E\big[(Y-r(\X))(r(\X)-g(\X))\big]
  &= \E\Big[\,\E\big[(Y-r(\X))(r(\X)-g(\X))\,\big|\,\X\big]\Big]\\
  &= \E\Big[\big(r(\X)-g(\X)\big)\,
       \underbrace{\E\big[Y-r(\X)\,\big|\,\X\big]}_{=\;r(\X)-r(\X)\;=\;0}\Big]
  = 0 .
\end{aligned}
\]
O passo do meio é a regra "o que é função de $\X$ sai da esperança condicional em
$\X$": $r(\X)-g(\X)$ depende só de $\X$, e condicionar em $\X$ é tratá-lo como
conhecido. O zero vem da \emph{definição} de $r$: $\E[Y\mid\X]=r(\X)$.

Em palavras: o erro $Y-r(\X)$ que sobra depois de acertar a média condicional é
\textbf{ortogonal a toda função de $\X$} --- nenhuma informação contida em $\X$ ajuda a
prevê-lo. Guarde a frase: o Exercício~\ref{ex:02-oraculo} tem uma versão dela para
as funções lineares.

\smallskip
\textbf{(b)} A identidade de (a) diz que
\[
  \risco(g) \;=\; \underbrace{\risco(r)}_{\text{não depende de } g}
   \;+\; \underbrace{\E\big[(r(\X)-g(\X))^2\big]}_{\ge\ 0} .
\]
Logo $\risco(g)\ge\risco(r)$ para toda $g$: $r$ é minimizadora. E é a única, a menos
de probabilidade zero: a igualdade vale se e só se $\E[(r(\X)-g(\X))^2]=0$, isto é,
se $g(\X)=r(\X)$ com probabilidade $1$. (Alterar $g$ num conjunto de valores de $\x$
que nunca ocorrem não muda risco nenhum.)

O valor mínimo sai de novo pela esperança total. Dado $\X$, $Y-r(\X)$ é o desvio de
$Y$ em torno da sua própria média condicional, e a esperança condicional do seu
quadrado é, por definição, a variância condicional:
\[
  \risco(r)=\E\big[(Y-\E[Y\mid\X])^2\big]
   =\E\Big[\,\E\big[(Y-\E[Y\mid\X])^2\;\big|\;\X\big]\Big]
   =\E\big[\Var(Y\mid\X)\big].
\]

\smallskip
\textbf{(c)} A impossibilidade é matemática, e não tecnológica. O item (b) mostrou que
\textbf{nenhuma} função de $\X$ tem risco menor que $\E[\Var(Y\mid\X)]$ --- nem mesmo
$r$, conhecida exatamente, com $n=\infty$ e erro de estimação zero. E essa quantidade
não menciona método nenhum: é uma propriedade da distribuição conjunta de $(\X,Y)$.
Ela mede quanto $Y$ ainda varia entre indivíduos com \emph{o mesmo} $\x$. Dois
pacientes com idade, peso e exames idênticos, mas desfechos diferentes, recebem
forçosamente a mesma previsão, porque o preditor recebeu a mesma entrada nos dois
casos.

Para baixar o número é preciso mudar o problema: \textbf{medir mais}. Se $Z$ é uma
covariável nova, a lei da variância total, aplicada condicionalmente a $\X$, dá
\[
  \E\big[\Var(Y\mid\X)\big]
  = \E\big[\Var(Y\mid\X,Z)\big]
    + \E\Big[\Var\big(\E[Y\mid\X,Z]\;\big|\;\X\big)\Big]
  \;\ge\; \E\big[\Var(Y\mid\X,Z)\big],
\]
com igualdade só quando $Z$ não muda a média condicional, $\E[Y\mid\X,Z]=\E[Y\mid\X]$.
Acrescentar informação nunca aumenta o erro irredutível, e o diminui sempre que a
informação for relevante. Isso vale na população; com $n$ finito, cada covariável
nova também custa variância de estimação (Aula~02), de modo que medir mais é
necessário, mas não suficiente, para prever melhor.

\smallskip
\textbf{(d)} Aplique a identidade de (a) com $g\equiv\E[Y]$. O lado esquerdo é o risco
do melhor preditor constante, $\Var(Y)$. E como $\E[r(\X)]=\E\big[\E[Y\mid\X]\big]=\E[Y]$,
pela esperança total, a última parcela é uma variância:
\[
  \Var(Y) = \risco(r) + \E\Big[\big(r(\X)-\E[r(\X)]\big)^2\Big]
  \qquad\Longrightarrow\qquad
  \boxed{\;\Var(Y)-\risco(r)=\Var\big(r(\X)\big)\;}
\]
É a lei da variância total, $\Var(Y)=\E[\Var(Y\mid\X)]+\Var(\E[Y\mid\X])$, obtida por
outro caminho.

A pergunta que isso responde é: \textbf{vale a pena usar as covariáveis?} O lado
esquerdo é o máximo que se pode ganhar trocando o preditor constante pelo melhor
preditor possível; o direito diz que esse ganho é exatamente o quanto a resposta
\emph{esperada} muda de um perfil de covariáveis para outro. Dividido por $\Var(Y)$,
vira a fração da variância de $Y$ que as covariáveis conseguem explicar no melhor dos
casos: um $R^2$ populacional, teto para o $R^2$ de qualquer modelo.

Os dois extremos ajudam a fixar. Se $\Var(r(\X))=0$, a função de regressão é
constante: as covariáveis não carregam informação sobre a média de $Y$, e prever pela
média é o melhor que existe --- não há modelo a construir. Se $\risco(r)=0$, $Y$ é
função de $\X$, e toda a sua variância é explicável.
\end{solucao}

\end{exercicio}

\begin{exercicio}[dois riscos com o mesmo nome, e uma barra de erro]
\label{ex:01-doisriscos}
% Aula 01 --- risco condicional x risco esperado, e a barra de erro do teste.
A palavra "risco" apareceu durante o nosso curso com dois sentidos, e a diferença entre eles
não é preciosismo: ela decide qual pergunta o seu número está respondendo.

Quando $g$ é tratada como uma função \emph{fixa}, a esperança em
$\E[(Y-g(\X))^2]$ percorre só a observação nova $(\X,Y)$. Chamamos isso de
\textbf{risco condicional}. Quando $g=\ghat$ é tratada como aleatória --- e ela é,
porque foi construída a partir de $\Dados$ ---, a esperança percorre também as
amostras de treino possíveis, e temos o \textbf{risco esperado}.
\begin{itens}
  \item \pts{1,3} Cada um dos dois oferece uma garantia diferente. Descreva as
        duas, e diga qual interessa a quem já ajustou um modelo e vai colocá-lo em
        produção, e qual interessa a quem ainda está escolhendo entre diversos métodos.
  \item \pts{1,2} Você separou um conjunto de teste de $m$ observações que não
        participou de nenhum treinamento, e calculou $\riscohat(\ghat)=\frac1m\sum_{k=1}^m W_k$,
        com $W_k=(Y_k-\ghat(\X_k))^2$. Qual dos dois riscos esse número estima?
        Para responder, pergunte-se o que está fixo e o que varia no momento em
        que você faz a conta.
  \item \pts{1,5} Um conjunto de teste permite algo que a validação cruzada não
        permite: uma barra de erro honesta. Explique por que os $W_k$ são i.i.d.\ e
        escreva o intervalo de confiança aproximado de 95\% que o Teorema Central
        do Limite fornece.
\end{itens}

\begin{solucao}
\textbf{(a)} O \textbf{risco condicional} é uma garantia sobre \emph{aquela
função específica}: se eu usar este $\ghat$, já ajustado, em observações novas, meu
erro quadrático médio esperado é tanto.

O \textbf{risco esperado} é uma garantia sobre o \emph{procedimento}: se eu
sortear uma amostra de treino, rodar este método e usar o que sair, meu erro
esperado é tanto. A esperança percorre também os treinos possíveis.

Quem vai pôr o modelo em produção quer o \textbf{condicional}. O modelo dele é um
só, já existe, não vai ser reajustado --- saber como se comportam os modelos que
aquele método \emph{costuma} produzir não paga as contas.

Quem está comparando métodos quer o \textbf{esperado}. A pergunta ali é "qual
deles tende a produzir bons modelos", e uma rodada de sorte com um método ruim não
deveria decidir a escolha.

\smallskip
\textbf{(b)} Estima o \textbf{risco condicional}.

No momento da conta, $\ghat$ já foi ajustado e não muda: é uma função fixa. O que
varia de uma parcela para outra é apenas o par $(\X_k,Y_k)$ sorteado para o
teste. A aleatoriedade do treino --- que outra amostra teria dado outro $\ghat$
--- não entra na conta em momento nenhum, porque o treino aconteceu antes e não é
repetido.

Dito de outro modo: para estimar o risco \emph{esperado} você teria de repetir o
processo inteiro, sorteando vários treinos, e não é o que um único conjunto de
teste faz.

\smallskip
\textbf{(c)} Condicionalmente a $\ghat$ --- que é a situação, por (b) --- cada
$W_k=(Y_k-\ghat(\X_k))^2$ é função de um único par $(\X_k,Y_k)$. Esses pares são
i.i.d.\ entre si e independentes do treino, e a mesma função $\ghat$ é aplicada a
todos. Logo os $W_k$ são i.i.d., com média $\risco(\ghat)$.

Como $\riscohat$ é média de i.i.d., o Teorema Central do Limite se aplica e dá
\[
  \riscohat(\ghat)\;\pm\;2\sqrt{\widehat\sigma^2/m},
  \qquad
  \widehat\sigma^2=\frac1m\sum_{k=1}^m(W_k-\overline W)^2 .
\]

Há uma hipótese fazendo todo o trabalho aqui: o teste não participou de nada. Se ele
tivesse sido usado para escolher qualquer coisa, os $W_k$ deixariam de ser
independentes do $\ghat$ e o intervalo perderia o sentido.
\end{solucao}

\end{exercicio}

%==============================================================================
\section*{Aula 02 --- Regressão Linear e Regularização}
%==============================================================================

\begin{exercicio}[de onde saem as fórmulas do caso ortonormal]
\label{ex:02-ortonormal}
% Aula 02 --- derivar as fórmulas do Ridge e do Lasso no caso ortonormal.
As notas da aula de Regressão Linear apresentam duas fórmulas prontas para o caso ortonormal,
$\mathbb{X}^\top\mathbb{X}=\mathbb{I}$, e não as derivam. São elas que explicam a
frase "o Ridge encolhe e o Lasso zera", e o exercício é preencher essa lacuna.

Use as definições da aula, \textbf{sem} fator $\tfrac12$ no RSS:
\[
  \widehat{\bbeta}^{\,\text{ridge}}
   = \argmin_{\bbeta}\Big\{\norm{\bm y-\mathbb{X}\bbeta}^2+\lambda\sum_j\beta_j^2\Big\},
  \qquad
  \widehat{\bbeta}^{\,\text{lasso}}
   = \argmin_{\bbeta}\Big\{\norm{\bm y-\mathbb{X}\bbeta}^2+\lambda\sum_j\abs{\beta_j}\Big\}.
\]
Aqui $\mathbb{X}\in\R^{n\times p}$ não tem a coluna de uns: não há intercepto, e as
duas penalidades somam sobre os $p$ coeficientes.
\begin{itens}
  \item \pts{1,1} Comece pelo passo que torna tudo possível. Supondo
        $\mathbb{X}^\top\mathbb{X}=\mathbb{I}$, mostre que o estimador de mínimos
        quadrados é $\widehat\bbeta=\mathbb{X}^\top\bm y$ e que
        \[
          \norm{\bm y-\mathbb{X}\bbeta}^2
            = \norm{\bm y}^2-\norm{\widehat\bbeta}^2
              + \sum_{j=1}^{p}\big(\beta_j-\widehat\beta_j\big)^2 .
        \]
        Explique por que essa identidade faz o problema \textbf{se separar} em $p$
        problemas de uma variável --- e por que é justamente isso que falha quando
        as colunas não são ortonormais.
  \item \pts{0,7} Com a separação em mãos, derive
        $\widehat\beta_j^{\,\text{ridge}}=\widehat\beta_j/(1+\lambda)$.
  \item \pts{1,4} Agora o Lasso. A dificuldade nova é que $\abs{\beta_j}$ não é
        derivável em zero, então não dá para igualar a derivada a zero e pronto:
        trate separadamente as regiões $\beta_j>0$, $\beta_j<0$ e o ponto
        $\beta_j=0$. Derive
        $\widehat\beta_j^{\,\text{lasso}}
          =\operatorname{sinal}(\widehat\beta_j)\big(\abs{\widehat\beta_j}-\lambda/2\big)_+$
        e diga qual é a condição exata que leva um coeficiente a zero.
  \item \pts{0,8} O limiar deu $\lambda/2$, e não $\lambda$. De onde veio esse
        $2$? Diga o que mudaria nas duas fórmulas se a aula tivesse definido o RSS
        com um $\tfrac12$ na frente, e que lição prática isso deixa.
\end{itens}

\begin{solucao}
\textbf{(a)} Com $\mathbb{X}^\top\mathbb{X}=\mathbb{I}$ a matriz é invertível, e as
equações normais $\mathbb{X}^\top\mathbb{X}\,\bbeta=\mathbb{X}^\top\bm y$ dão
diretamente $\widehat\bbeta=\mathbb{X}^\top\bm y$.

Para a identidade, o caminho mais curto é o teorema de Pitágoras. Escreva
\[
  \bm y-\mathbb{X}\bbeta=(\bm y-\mathbb{X}\widehat\bbeta)+\mathbb{X}(\widehat\bbeta-\bbeta).
\]
As duas parcelas são ortogonais, porque as equações normais dizem exatamente que o
resíduo do MQO é ortogonal às colunas de $\mathbb{X}$:
$\mathbb{X}^\top(\bm y-\mathbb{X}\widehat\bbeta)=\bm 0$. Logo, para \emph{qualquer}
$\mathbb{X}$ de posto cheio,
\[
  \norm{\bm y-\mathbb{X}\bbeta}^2
  =\norm{\bm y-\mathbb{X}\widehat\bbeta}^2
   +(\bbeta-\widehat\bbeta)^\top\mathbb{X}^\top\mathbb{X}\,(\bbeta-\widehat\bbeta).
  \tag{$*$}
\]
No caso ortonormal a forma quadrática é $\sum_j(\beta_j-\widehat\beta_j)^2$, e
\[
  \norm{\bm y-\mathbb{X}\widehat\bbeta}^2
  =\norm{\bm y}^2-2\,\widehat\bbeta^\top\underbrace{\mathbb{X}^\top\bm y}_{=\;\widehat\bbeta}
   +\widehat\bbeta^\top\underbrace{\mathbb{X}^\top\mathbb{X}}_{=\;\mathbb{I}}\widehat\bbeta
  =\norm{\bm y}^2-\norm{\widehat\bbeta}^2 ,
\]
o que dá a identidade do enunciado.

\emph{Por que ela separa o problema.} Somando a penalidade, os dois objetivos ficam
\[
  \norm{\bm y}^2-\norm{\widehat\bbeta}^2+\sum_{j=1}^p h_j(\beta_j),
  \qquad
  h_j(\beta)=(\beta-\widehat\beta_j)^2+\lambda\,\mathrm{pen}(\beta),
\]
com $\mathrm{pen}(\beta)=\beta^2$ no Ridge e $\mathrm{pen}(\beta)=\abs{\beta}$ no
Lasso. A constante não muda o minimizador, e o resto é uma soma de funções de
\emph{uma coordenada cada}. Para somas assim vale um fato simples: $\bbeta^\ast$
minimiza $\sum_jh_j(\beta_j)$ se e só se cada $\beta_j^\ast$ minimiza a sua $h_j$. (Se
cada parcela está no seu mínimo, a soma está; e se alguma não estivesse, trocar só
aquela coordenada baixaria a soma.) Um problema em $\R^p$ virou $p$ problemas na reta.

\emph{Por que isso falha sem ortonormalidade.} Em $(*)$, a forma quadrática tem os
termos cruzados $2(\mathbb{X}^\top\mathbb{X})_{j\ell}(\beta_j-\widehat\beta_j)(\beta_\ell-\widehat\beta_\ell)$,
e o melhor valor de uma coordenada passa a depender das outras. É por isso que o Lasso
não tem fórmula fechada no caso geral, e as fórmulas deste exercício pedem uma
hipótese forte. (É também por isso que o enunciado tira o intercepto: a coluna de uns
tem norma $\sqrt n$ e não caberia numa matriz com $\mathbb{X}^\top\mathbb{X}=\mathbb{I}$.
Com dados centrados o intercepto sai de cena, e é nesse cenário que o caso ortonormal
faz sentido.)

\smallskip
\textbf{(b)} Para o Ridge, $h(\beta)=(\beta-\widehat\beta_j)^2+\lambda\beta^2$ é uma
parábola com $h''=2(1+\lambda)>0$: estritamente convexa, com mínimo único onde a
derivada se anula,
\[
  h'(\beta)=2(\beta-\widehat\beta_j)+2\lambda\beta=0
  \iff
  \widehat\beta_j^{\,\text{ridge}}=\frac{\widehat\beta_j}{1+\lambda}.
\]
O Ridge \textbf{multiplica} todo coeficiente pelo mesmo fator $1/(1+\lambda)$: encolhe
proporcionalmente --- quem é grande perde mais em valor absoluto --- e nunca zera um
coeficiente não nulo, para $\lambda$ finito.

\smallskip
\textbf{(c)} Agora $h(\beta)=(\beta-b)^2+\lambda\abs{\beta}$, com $b=\widehat\beta_j$ e
$\lambda>0$. A função é contínua em toda a reta e derivável fora do zero:
\[
  h'(\beta)=2(\beta-b)+\lambda\quad(\beta>0),
  \qquad
  h'(\beta)=2(\beta-b)-\lambda\quad(\beta<0).
\]
A estratégia é ver onde $h$ cresce e onde decresce, dos dois lados do zero.

\emph{Caso $b>\lambda/2$.} Para $\beta<0$, $h'(\beta)=2\beta-(2b+\lambda)<0$: $h$
decresce em toda a semirreta negativa. Para $\beta>0$, $h'$ troca de sinal uma única
vez, em $\beta=b-\lambda/2>0$: $h$ decresce até ali e cresce depois. Como $h$ é
contínua no zero, o mínimo global é $\beta=b-\lambda/2$.

\emph{Caso $b<-\lambda/2$.} Simétrico: o mínimo é $\beta=b+\lambda/2<0$.

\emph{Caso $\abs{b}\le\lambda/2$.} Para $\beta>0$, $h'(\beta)=2\beta+(\lambda-2b)>0$,
porque $\lambda-2b\ge0$; para $\beta<0$, $h'(\beta)=2\beta-(\lambda+2b)<0$, porque
$\lambda+2b\ge0$. A função decresce até o zero e cresce depois dele: o mínimo é
$\beta=0$ --- justamente o ponto em que ela não é derivável, e que nenhuma equação
"$h'=0$" encontraria.

Os três casos cabem numa fórmula só,
\[
  \widehat\beta_j^{\,\text{lasso}}
  =\operatorname{sinal}(\widehat\beta_j)\big(\abs{\widehat\beta_j}-\lambda/2\big)_+ ,
  \qquad\text{e o coeficiente zera}\iff
  \boxed{\;\abs{\widehat\beta_j}\le\lambda/2\;}.
\]
(Para quem conhece subgradientes: $0$ minimiza a função convexa $h$ quando
$0\in\partial h(0)=-2b+\lambda[-1,1]$, que é a mesma condição.)

Compare com (b): o Lasso \textbf{subtrai} a mesma quantidade, $\lambda/2$, de todo
coeficiente, em vez de multiplicar por um fator. Quem tinha módulo menor que isso não
sobrevive ao deslocamento e vira exatamente zero. É o \emph{soft thresholding}, e é a
razão algébrica de o Lasso fazer seleção de variáveis.

\smallskip
\textbf{(d)} O $2$ vem de derivar o quadrado. A derivada de $(\beta-\widehat\beta_j)^2$
traz um fator $2$ que a de $\lambda\abs{\beta}$ não traz; ao isolar $\beta$, dividimos
tudo por $2$, e o $\lambda$ vira $\lambda/2$.

Com o RSS definido com $\tfrac12$ na frente (e a penalidade sem), o problema seria
minimizar $\tfrac12\norm{\bm y-\mathbb{X}\bbeta}^2+\lambda\,\mathrm{Pen}(\bbeta)$.
Multiplicar um objetivo por uma constante positiva não muda o minimizador, e
multiplicando este por $2$ ele vira o problema da aula com $2\lambda$ no lugar de
$\lambda$. Basta fazer essa troca nas fórmulas:
\[
  \widehat\beta_j^{\,\text{ridge}}=\frac{\widehat\beta_j}{1+2\lambda},
  \qquad
  \widehat\beta_j^{\,\text{lasso}}
  =\operatorname{sinal}(\widehat\beta_j)\big(\abs{\widehat\beta_j}-\lambda\big)_+ .
\]
(Pôr o $\tfrac12$ também na penalidade não mudaria nada: seria multiplicar o objetivo
inteiro por uma constante.)

A lição prática: \textbf{o valor numérico de $\lambda$ depende da convenção}, não só
dos dados. Antes de comparar o $\lambda$ das notas com o de um livro, de um artigo ou
de uma biblioteca, confira qual objetivo cada um minimiza. O \texttt{Lasso} do
\texttt{scikit-learn} minimiza $\tfrac1{2n}\norm{\bm y-\mathbb{X}\bbeta}^2+\alpha\norm{\bbeta}_1$;
multiplicando por $2n$ vê-se que o seu $\alpha$ corresponde a $\lambda=2n\alpha$ na
convenção da aula. O $n$ vem de usar a média em vez da soma; o $2$ é o deste item.
\end{solucao}

\end{exercicio}

\begin{exercicio}[o que o MQO estima quando a linearidade é falsa]
\label{ex:02-oraculo}
% Aula 02 --- o melhor preditor linear (o "oráculo") quando r não é linear.
Circula por aí a ideia de que a regressão linear "só vale se o modelo estiver
correto". Vimos durante o curso que isso é impreciso: o MQO converge para uma função
bem definida mesmo quando $r$ não é linear.

Essa função é o \textbf{melhor preditor linear},
\[
  \bbeta^\ast = \argmin_{\bbeta} \E\big[(Y-\bbeta^\top\X)^2\big],
\]
que existe sempre que $\Sigma=\E[\X\X^\top]$ é bem definida e invertível ---
\emph{sem} supor nada sobre a forma de $r$.
\begin{itens}
  \item \pts{0,9} Derive $\bbeta^\ast=\Sigma^{-1}\bm\alpha$, onde
        $\bm\alpha=\E[Y\X]$. É o mesmo cálculo das equações normais, com
        esperanças no lugar de somas.
  \item \pts{1,0} O estimador de mínimos quadrados pode ser escrito como
        $\widehat\bbeta=\widehat\Sigma^{-1}\widehat{\bm\alpha}$, com
        $\widehat\Sigma=\frac1n\sum_i\X_i\X_i^\top$ e
        $\widehat{\bm\alpha}=\frac1n\sum_i Y_i\X_i$. Olhe para essas duas
        expressões e diga o que elas são, estatisticamente. Em seguida, explique
        por que $\widehat\bbeta\to\bbeta^\ast$ em probabilidade sem que se precise
        supor linearidade em lugar nenhum.
  \item \pts{1,0} Todo teorema garante uma coisa e deixa de garantir outra.
        Enuncie as duas, com cuidado. Em particular: o resultado diz que o risco
        de $\widehat\bbeta$ é pequeno?
  \item \pts{1,1} Um colega ajusta uma regressão linear a dados em que $r$ é
        visivelmente curva e conclui: "o coeficiente deu $2{,}3$, logo aumentar
        $x_1$ em uma unidade aumenta $Y$ em $2{,}3$". Separe o que é leitura
        legítima de $\bbeta^\ast$ do que não é.
\end{itens}

\begin{solucao}
Hipóteses em uso: $\E[Y^2]<\infty$ e $\E\norm{\X}^2<\infty$, para que $\Sigma$ e
$\bm\alpha$ existam (a existência de $\bm\alpha$ sai por Cauchy--Schwarz), e $\Sigma$
invertível. Como $\bm v^\top\Sigma\bm v=\E\big[(\bm v^\top\X)^2\big]\ge0$ para todo
$\bm v$, $\Sigma$ é semidefinida positiva; sendo invertível, é definida positiva.

\smallskip
\textbf{(a)} Expanda o objetivo; as esperanças são números fixos:
\[
  Q(\bbeta)=\E\big[(Y-\bbeta^\top\X)^2\big]
  =\E[Y^2]-2\,\bbeta^\top\underbrace{\E[Y\X]}_{\bm\alpha}
   +\bbeta^\top\underbrace{\E[\X\X^\top]}_{\Sigma}\bbeta .
\]
É uma função quadrática de $\bbeta$, com gradiente $\nabla Q(\bbeta)=-2\bm\alpha+2\Sigma\bbeta$
e Hessiana $2\Sigma$, definida positiva. Logo $Q$ é estritamente convexa, e seu único
ponto crítico é o mínimo global:
\[
  \Sigma\bbeta=\bm\alpha
  \iff
  \bbeta^\ast=\Sigma^{-1}\bm\alpha .
\]
São as equações normais da aula com $\E[\X\X^\top]$ no lugar de
$\frac1n\mathbb{X}^\top\mathbb{X}$ e $\E[Y\X]$ no lugar de $\frac1n\mathbb{X}^\top\bm y$:
o mesmo problema, escrito na população em vez de na amostra.

Vale conhecer uma segunda demonstração, sem derivadas, porque ela rende uma identidade
que usaremos em (c). Some e subtraia $\bbeta^{\ast\top}\X$ dentro do quadrado:
\[
\begin{aligned}
  Q(\bbeta)
  &=\E\Big[\big((Y-\bbeta^{\ast\top}\X)-(\bbeta-\bbeta^\ast)^\top\X\big)^2\Big]\\
  &=Q(\bbeta^\ast)
   -2(\bbeta-\bbeta^\ast)^\top\underbrace{\E\big[\X\,(Y-\X^\top\bbeta^\ast)\big]}_{=\;\bm\alpha-\Sigma\bbeta^\ast\;=\;\bm 0}
   +(\bbeta-\bbeta^\ast)^\top\Sigma\,(\bbeta-\bbeta^\ast),
\end{aligned}
\]
ou seja,
\[
  Q(\bbeta)=Q(\bbeta^\ast)+(\bbeta-\bbeta^\ast)^\top\Sigma\,(\bbeta-\bbeta^\ast).
  \tag{$\dagger$}
\]
O último termo é estritamente positivo para $\bbeta\neq\bbeta^\ast$, o que prova de uma
vez que $\bbeta^\ast$ é o mínimo e que é o único. É o mesmo "some e subtraia" do
Exercício~\ref{ex:01-ropt}, e funciona pela mesma razão. Lá, o resíduo $Y-r(\X)$ era
ortogonal a toda função de $\X$; aqui, $\E\big[\X(Y-\X^\top\bbeta^\ast)\big]=\bm0$ diz
que o resíduo do melhor preditor linear é ortogonal a cada covariável. $r(\X)$ é a
projeção de $Y$ sobre todas as funções de $\X$; $\bbeta^{\ast\top}\X$, a projeção sobre
as lineares.

\smallskip
\textbf{(b)} $\widehat\Sigma$ e $\widehat{\bm\alpha}$ são \textbf{médias amostrais} de
variáveis i.i.d.: a primeira, das matrizes $\X_i\X_i^\top$, cuja esperança é $\Sigma$;
a segunda, dos vetores $Y_i\X_i$, cuja esperança é $\bm\alpha$. É essa observação que
resolve o item.

Pela lei dos grandes números, aplicada coordenada a coordenada (as esperanças existem
pelas hipóteses acima), $\widehat\Sigma\to\Sigma$ e $\widehat{\bm\alpha}\to\bm\alpha$ em
probabilidade. A função $(A,\bm a)\mapsto A^{-1}\bm a$ é contínua no conjunto aberto das
matrizes invertíveis, que contém $\Sigma$; por isso $\widehat\Sigma$ é invertível com
probabilidade tendendo a $1$, e o teorema de Mann--Wald (da aplicação contínua)
transporta a convergência:
\[
  \widehat\bbeta=\widehat\Sigma^{-1}\widehat{\bm\alpha}
  \;\xrightarrow{\ P\ }\;\Sigma^{-1}\bm\alpha=\bbeta^\ast .
\]
Em nenhum passo apareceu a forma de $r$. O alvo $\bbeta^\ast$ é definido por um
problema de minimização que faz sentido para qualquer distribuição com segundos
momentos finitos; o que a linearidade de $r$ acrescentaria seria só a coincidência
$\bbeta^{\ast\top}\x=r(\x)$ --- agradável, mas dispensável para a convergência.

\smallskip
\textbf{(c)} Duas identidades do curso dizem tudo. Aplique a do
Exercício~\ref{ex:01-ropt}(a) com $g(\x)=\bbeta^{\ast\top}\x$, e depois $(\dagger)$ com
$\bbeta=\widehat\bbeta$ (fixado o treino, o risco do ajuste é
$\risco(\widehat g)=Q(\widehat\bbeta)$):
\[
  \risco(\widehat g)=
  \underbrace{\E\big[\Var(Y\mid\X)\big]}_{\text{irredutível}}
  +\underbrace{\E\big[(r(\X)-\bbeta^{\ast\top}\X)^2\big]}_{\text{erro de aproximação}}
  +\underbrace{(\widehat\bbeta-\bbeta^\ast)^\top\Sigma\,(\widehat\bbeta-\bbeta^\ast)}_{\text{erro de estimação}} .
\]
Com a decomposição à vista, as duas coisas se enunciam sem ambiguidade.

\textbf{O que o teorema garante:} a terceira parcela vai a zero, porque
$\widehat\bbeta\to\bbeta^\ast$ e a forma quadrática é contínua. Com $n$ grande, o MQO
chega tão perto quanto se queira do melhor que a classe das funções lineares tem a
oferecer, e $\risco(\widehat g)\to\risco(g_{\bbeta^\ast})$.

\textbf{O que ele não garante:} que esse melhor seja bom. A segunda parcela não depende
de $n$, e o teorema não diz nada sobre ela; se $r$ é muito curva, ela é grande, e
convergir para uma aproximação ruim continua sendo ruim. Então não: o resultado
\emph{não} diz que o risco de $\widehat\bbeta$ é pequeno, só que ele se aproxima do
risco do oráculo. Tampouco diz \emph{com que rapidez}: é um resultado de convergência,
sem taxa.

Na linguagem da Aula~01, o teorema cuida do erro de estimação --- a variância, e a
parte do viés que vem de $n$ ser finito --- e deixa intacto o viés de aproximação, que
é o preço fixo de ter escolhido uma classe rígida.

\smallskip
\textbf{(d)} \emph{A leitura legítima.} $2{,}3$ é o coeficiente de $x_1$ no hiperplano
que \emph{melhor aproxima} $r$ em erro quadrático médio, sob a distribuição de $\X$ em
que os dados foram colhidos. É um resumo global, e útil como tal.

\emph{O que não se pode ler nele}, em três sentidos.
\begin{itemize}[leftmargin=1.2cm]
  \item \emph{Não é a inclinação de $r$.} Com $r$ curva, $\partial r/\partial x_1$ muda
        de ponto para ponto, e $2{,}3$ não é o valor dela em lugar nenhum em
        particular. Com várias covariáveis, pode nem ter relação com ela. Tome
        $X_2\sim N(0,1)$, $X_1=X_2^2+\varepsilon$, com $\varepsilon\sim N(0,s^2)$
        independente, e $r(\x)=x_2^2$: $r$ nem depende de $x_1$, e ainda assim,
        regredindo em $(1,x_1,x_2)$,
        \[
          \beta_1^\ast=\frac{\Cov(X_1,X_2^2)}{\Var(X_1)}=\frac{2}{2+s^2}\neq0 .
        \]
        A reta usa $x_1$ como \emph{procurador} da parte curva de $r$ que ela não
        consegue representar em $x_2$. (A conta: $X_2$ não é correlacionada nem com
        $X_2^2$ nem com $X_1$, pois $\E[X_2^3]=0$; então o coeficiente de $x_2$ é zero,
        o de $x_1$ é o da regressão simples, e $\Var(X_2^2)=\E[X_2^4]-1=2$.)
  \item \emph{Depende de onde os dados caíram.} Mesmo com uma covariável só: para
        $r(x)=x^2$ e $X\sim\mathrm{Unif}(0,a)$, a inclinação ótima é
        $\Cov(X,X^2)/\Var(X)=(a^3/12)/(a^2/12)=a$. Colhendo dados em $[0,1]$ o colega
        veria $1$; em $[0,2]$, veria $2$ --- com $r$ exatamente a mesma. Um número que
        muda com a faixa amostrada não descreve um mecanismo.
  \item \emph{"Aumentar $x_1$ aumenta $Y$" é uma frase causal}, sobre intervir no
        sistema. Nenhuma regressão a autoriza sozinha, nem com $r$ linear: comparar
        indivíduos que \emph{diferem} em $x_1$ não é o mesmo que \emph{mudar} o $x_1$
        de alguém.
\end{itemize}
\end{solucao}

\end{exercicio}

\begin{exercicio}[a Ridge como moda a posteriori]
\label{ex:02-bayes}
% Aula 02 --- a Ridge como moda a posteriori sob priori gaussiana.
Já mencionamos que o Ridge tem uma leitura bayesiana: ele é a moda a
posteriori sob uma priori gaussiana. A afirmação é mais do que curiosidade ---
ela dá um significado concreto ao $\lambda$, que até aqui era só um hiperparâmetro.

Considere o modelo
\[
  \bm y \mid \bbeta \;\sim\; N\big(\mathbb{X}\bbeta,\ \sigma^2\mathbb{I}\big),
  \qquad
  \bbeta \;\sim\; N\big(\bm 0,\ \sigma_\beta^2\mathbb{I}\big),
\]
com $\sigma^2$ e $\sigma_\beta^2$ conhecidos.
\begin{itens}
  \item \pts{1,6} Escreva a densidade a posteriori de $\bbeta$ a menos de
        constante multiplicativa e mostre que maximizá-la equivale a minimizar
        \[
          \frac{1}{2\sigma^2}\norm{\bm y-\mathbb{X}\bbeta}^2
          + \frac{1}{2\sigma_\beta^2}\norm{\bbeta}^2 .
        \]
        O caminho é o de sempre: passar ao logaritmo transforma produto em soma, e
        trocar o sinal transforma maximizar em minimizar.
  \item \pts{1,0} Compare o que você obteve com a definição do Ridge da aula e
        identifique $\lambda$ em função de $\sigma^2$ e $\sigma_\beta^2$. Cuidado
        com o fator que precisa ser eliminado antes da comparação.
  \item \pts{1,4} Agora a parte interessante: o que a priori está \emph{afirmando}
        quando $\sigma_\beta^2\to\infty$, e quando $\sigma_\beta^2\to 0$? Ligue
        cada resposta ao comportamento do Ridge em $\lambda\to 0$ e
        $\lambda\to\infty$.
\end{itens}

\begin{solucao}
\textbf{(a)} Pelo teorema de Bayes, a posteriori é proporcional ao produto da
verossimilhança pela priori:
\[
  \pi(\bbeta\mid\bm y)\;\propto\; f(\bm y\mid\bbeta)\,\pi(\bbeta)
   \;\propto\;
   \exp\Big\{-\frac{1}{2\sigma^2}\norm{\bm y-\mathbb{X}\bbeta}^2\Big\}
   \cdot
   \exp\Big\{-\frac{1}{2\sigma_\beta^2}\norm{\bbeta}^2\Big\}.
\]
As constantes de normalização das duas gaussianas não dependem de $\bbeta$ e por
isso podem ser absorvidas no $\propto$.

Maximizar uma função positiva é o mesmo que maximizar seu logaritmo, e o
logaritmo de um produto de exponenciais é a soma dos expoentes:
\[
  \log\pi(\bbeta\mid\bm y)
   = \text{const} - \frac{1}{2\sigma^2}\norm{\bm y-\mathbb{X}\bbeta}^2
     - \frac{1}{2\sigma_\beta^2}\norm{\bbeta}^2 .
\]
Maximizar isso é minimizar o simétrico dos dois termos que dependem de $\bbeta$,
que é exatamente o objetivo do enunciado.

\smallskip
\textbf{(b)} O objetivo de (a) tem um $1/(2\sigma^2)$ na frente do primeiro termo,
e a definição do Ridge não tem fator nenhum ali. Multiplique tudo por
$2\sigma^2>0$ --- o que não altera o minimizador, por ser constante positiva:
\[
  \norm{\bm y-\mathbb{X}\bbeta}^2 + \frac{\sigma^2}{\sigma_\beta^2}\norm{\bbeta}^2 .
\]
Agora a comparação é imediata, e
\[
  \boxed{\;\lambda=\frac{\sigma^2}{\sigma_\beta^2}\;}
\]

\smallskip
\textbf{(c)} Quando $\sigma_\beta^2\to\infty$, a priori fica \textbf{difusa}: ela
espalha massa por toda parte e não expressa preferência alguma por coeficientes
pequenos. Em outras palavras, ela não afirma nada. Correspondentemente
$\lambda\to 0$ e o Ridge degenera no MQO: os dados decidem sozinhos, com viés
baixo e variância alta.

Quando $\sigma_\beta^2\to 0$, a priori fica \textbf{concentrada em zero}: ela
afirma, antes de ver dado nenhum, que os coeficientes são praticamente nulos.
Correspondentemente $\lambda\to\infty$ e $\widehat\bbeta\to\bm 0$: viés máximo,
variância mínima.

Entre os dois extremos, a fórmula $\lambda=\sigma^2/\sigma_\beta^2$ diz uma coisa
bem razoável: regulariza-se mais quando os \emph{dados} são ruidosos
($\sigma^2$ grande) ou quando a \emph{crença prévia} é firme ($\sigma_\beta^2$
pequeno). O $\lambda$ é a razão entre duas incertezas, e escolhê-lo por validação
cruzada é, nessa leitura, deixar os dados escolherem a priori que melhor prediz.
\end{solucao}

\end{exercicio}

%==============================================================================
\section*{Aula 03 --- Seleção de Modelos e Validação Cruzada}
%==============================================================================

\begin{exercicio}[o atalho do LOOCV, demonstrado]
\label{ex:03-atalho}
% Aula 03 --- demonstrar o atalho do LOOCV para ajustes lineares.
O LOOCV (\emph{Leave-One-Out Cross Validation}) parece caro: $n$ ajustes, um para cada observação deixada de fora.
Já comentamos que, para ajustes lineares, um único ajuste basta --- mas não demonstramos.

Vamos demonstrar no caso da regressão linear por mínimos quadrados. A matriz de
delineamento $\mathbb{X}$ tem linha $i$ igual a $\X_i^\top$, e o ajuste com
\emph{todos} os dados é $\widehat g(\x)=\x^\top\widehat\bbeta$, com
$\widehat\bbeta=(\mathbb{X}^\top\mathbb{X})^{-1}\mathbb{X}^\top\bm y$. Chame de
$\widehat g_{-i}$ o ajuste feito sem a $i$-ésima observação. A fórmula a demonstrar
é
\[
  \riscohat_{\text{LOO}}
   = \frac1n\sum_{i=1}^{n}\left(\frac{Y_i-\widehat g(\X_i)}{1-h_{ii}}\right)^{2},
\]
em que $h_{ii}$ é o $i$-ésimo elemento da diagonal da matriz
$\bm H=\mathbb{X}(\mathbb{X}^\top\mathbb{X})^{-1}\mathbb{X}^\top$.

Suponha que $\mathbb{X}$ tenha posto cheio, e que continue tendo quando se tira
qualquer uma de suas linhas. Assim todos os ajustes são únicos, e $h_{ii}\neq 1$:
dividir por $1-h_{ii}$ nunca é problema.

A demonstração vem em etapas. O item (a) mostra o que são os números $h_{i\ell}$;
os itens (b) e (c) usam um truque para escrever $\widehat g_{-i}(\X_i)$ em termos de
$\widehat g$; o (d) faz a conta; e o (e) interpreta o resultado.
\begin{itens}
  \item \pts{0,6} \emph{De onde vêm os pesos.} Mostre que o vetor de valores
        ajustados, $\widehat{\bm y}=\big(\widehat g(\X_1),\dots,\widehat g(\X_n)\big)^\top$,
        é igual a $\bm H\bm y$, e conclua, olhando a coordenada $i$, que
        \[
          \widehat g(\X_i)=\sum_{\ell=1}^{n} h_{i\ell}\,Y_\ell ,
        \]
        em que $h_{i\ell}$ é a entrada $(i,\ell)$ de $\bm H$. Dica: como a linha
        $i$ de $\mathbb{X}$ é $\X_i^\top$, vale $\widehat{\bm y}=\mathbb{X}\widehat\bbeta$.

        Em palavras: a predição no ponto $\X_i$ é uma combinação de todas as
        respostas, inclusive da própria $Y_i$, que entra com peso $h_{ii}$. De que
        esses pesos dependem? E de que eles \emph{não} dependem?
  \item \pts{1,0} \emph{O truque.} Imagine o conjunto de dados em que a resposta
        $Y_i$ foi trocada por $\widetilde Y_i=\widehat g_{-i}(\X_i)$ --- o valor que o
        ajuste sem a observação $i$ prevê ali ---, com as covariáveis e as demais
        respostas intactas. Mostre que o ajuste de mínimos quadrados nesse conjunto
        modificado é exatamente $\widehat g_{-i}$.

        Antes do caso geral, veja o truque funcionar no modelo só com intercepto,
        em que o ajuste é a média: troque $Y_i$ pela média das outras $n-1$
        respostas e confira que a média das $n$ passa a ser justamente a média das
        outras. Para o caso geral, um caminho: escreva o RSS do conjunto modificado
        como a soma das $n-1$ parcelas de sempre com a parcela da observação $i$;
        veja o que $\widehat g_{-i}$ faz com cada uma das duas partes; e conclua
        que nenhuma função linear consegue um RSS menor, o que basta, porque o
        ajuste é único.
  \item \pts{0,8} \emph{A mesma matriz, duas vezes.} O conjunto modificado tem as
        mesmas covariáveis que o original, e portanto a mesma matriz $\bm H$.
        Aplique (a) a ele e use (b) para mostrar que
        \[
          \widehat g_{-i}(\X_i)=\sum_{\ell\neq i} h_{i\ell}\,Y_\ell + h_{ii}\,\widetilde Y_i .
        \]
        Depois some e subtraia $h_{ii}Y_i$ para chegar a
        \[
          \widehat g_{-i}(\X_i) = \widehat g(\X_i) + h_{ii}\big(\widetilde Y_i - Y_i\big).
        \]
  \item \pts{0,7} \emph{A conta.} Na igualdade de (c), $\widetilde Y_i$ aparece dos
        dois lados, porque o lado esquerdo é o próprio $\widetilde Y_i$. Subtraia
        $Y_i$ dos dois lados, junte os termos com $\widetilde Y_i-Y_i$ e isole o
        resíduo de validação $Y_i-\widehat g_{-i}(\X_i)$. Conclua a fórmula do
        enunciado.
  \item \pts{0,9} \emph{Lendo o resultado.} A matriz $\bm H$ tem duas propriedades.
        É simétrica, porque $\mathbb{X}^\top\mathbb{X}$ é simétrica e sua inversa
        também. E vale $\bm H\bm H=\bm H$, o que se confere multiplicando:
        \[
          \bm H\bm H
           = \mathbb{X}(\mathbb{X}^\top\mathbb{X})^{-1}
             \underbrace{\mathbb{X}^\top\mathbb{X}(\mathbb{X}^\top\mathbb{X})^{-1}}_{=\;\mathbb{I}}
             \mathbb{X}^\top
           = \bm H .
        \]
        Em palavras: $\bm H$ projeta $\bm y$ no espaço gerado pelas colunas de
        $\mathbb{X}$, e projetar de novo o que já foi projetado não muda nada.
        Usando as duas propriedades, a entrada $(i,i)$ de $\bm H\bm H=\bm H$ dá
        \[
          h_{ii}=\sum_{\ell=1}^{n} h_{i\ell}^2 = h_{ii}^2+\sum_{\ell\neq i} h_{i\ell}^2 .
        \]
        Conclua que $0\le h_{ii}\le 1$ e que, portanto, o resíduo de validação
        nunca é menor, em módulo, que o de treino. O que isso confirma sobre o erro
        de treino como estimativa do risco? E o que significa uma observação com
        $h_{ii}$ próximo de $1$?
\end{itens}

\begin{solucao}
\textbf{(a)} Como a linha $i$ de $\mathbb{X}$ é $\X_i^\top$, a coordenada $i$ de
$\mathbb{X}\widehat\bbeta$ é $\X_i^\top\widehat\bbeta=\widehat g(\X_i)$; logo
$\widehat{\bm y}=\mathbb{X}\widehat\bbeta$ e, substituindo $\widehat\bbeta$,
\[
  \widehat{\bm y}
  =\underbrace{\mathbb{X}(\mathbb{X}^\top\mathbb{X})^{-1}\mathbb{X}^\top}_{=\;\bm H}\,\bm y
  =\bm H\bm y .
\]
A coordenada $i$ de $\bm H\bm y$ é a linha $i$ de $\bm H$ vezes $\bm y$; igualando as
coordenadas $i$ dos dois lados, $\widehat g(\X_i)=\sum_\ell h_{i\ell}Y_\ell$.

Os pesos são entradas de $\bm H$, que é construída só com $\mathbb{X}$: eles
\textbf{dependem das covariáveis e não dependem das respostas}. Troque o vetor $\bm y$
mantendo as covariáveis, e a predição em $\X_i$ continua sendo a mesma combinação,
agora das respostas novas. É essa propriedade --- o ajuste é uma função \emph{linear}
de $\bm y$ --- que o item (c) usa. (No modelo só com intercepto, $\mathbb{X}=\bm1$, todas
as entradas de $\bm H$ valem $1/n$, e a combinação é a média.)

\smallskip
\textbf{(b)} \emph{O caso da média.} Seja $\bar Y_{-i}$ a média das outras $n-1$
respostas. Trocando $Y_i$ por $\widetilde Y_i=\bar Y_{-i}$, a soma das $n$ respostas vira
$(n-1)\bar Y_{-i}+\bar Y_{-i}=n\bar Y_{-i}$, e a nova média é $\bar Y_{-i}$. A observação
trocada foi posta exatamente onde o ajuste sem ela já estava, e por isso não puxa o
ajuste para lado nenhum.

\emph{O caso geral} é o mesmo argumento, escrito com o RSS. Para uma função linear
$g(\x)=\x^\top\bbeta$ qualquer, o RSS do conjunto modificado se parte em duas:
\[
  \widetilde{\mathrm{RSS}}(g)
  =\underbrace{\sum_{\ell\neq i}\big(Y_\ell-g(\X_\ell)\big)^2}_{A(g)}
  +\underbrace{\big(\widetilde Y_i-g(\X_i)\big)^2}_{B(g)\;\ge\;0} .
\]
Por definição, $\widehat g_{-i}$ minimiza $A$ entre as funções lineares --- e está bem
definido, porque $\mathbb{X}$ sem a linha $i$ tem posto cheio. E $B(\widehat g_{-i})=0$,
porque $\widetilde Y_i$ foi escolhido como $\widehat g_{-i}(\X_i)$. Logo, para toda $g$
linear,
\[
  \widetilde{\mathrm{RSS}}(g)=A(g)+B(g)\;\ge\;A(\widehat g_{-i})+0
  =\widetilde{\mathrm{RSS}}(\widehat g_{-i}):
\]
$\widehat g_{-i}$ minimiza o RSS modificado. Como o conjunto modificado tem a mesma
matriz $\mathbb{X}$, de posto cheio, esse minimizador é único; é, portanto, o ajuste de
mínimos quadrados do conjunto modificado.

\smallskip
\textbf{(c)} Aplique (a) ao conjunto modificado. As covariáveis são as mesmas, então
$\bm H$ é a mesma, e a predição em $\X_i$ é a mesma combinação de antes, com
$\widetilde Y_i$ no lugar de $Y_i$. Por (b), essa predição é $\widehat g_{-i}(\X_i)$:
\[
  \widehat g_{-i}(\X_i)=\sum_{\ell\neq i}h_{i\ell}Y_\ell+h_{ii}\widetilde Y_i .
\]
Somando e subtraindo $h_{ii}Y_i$, a soma se completa e aparece $\widehat g(\X_i)$:
\[
  \widehat g_{-i}(\X_i)
  =\underbrace{\sum_{\ell}h_{i\ell}Y_\ell}_{=\;\widehat g(\X_i)}+h_{ii}\big(\widetilde Y_i-Y_i\big).
\]

\smallskip
\textbf{(d)} O lado esquerdo de (c) é o próprio $\widetilde Y_i$. Subtraindo $Y_i$ dos
dois lados, a diferença $\widetilde Y_i-Y_i$ aparece nos dois:
\[
  \widetilde Y_i-Y_i=\big(\widehat g(\X_i)-Y_i\big)+h_{ii}\big(\widetilde Y_i-Y_i\big)
  \iff
  (1-h_{ii})\big(\widetilde Y_i-Y_i\big)=\widehat g(\X_i)-Y_i .
\]
Como $h_{ii}\neq1$ (o item (e) confere), podemos dividir; trocando os sinais,
\[
  \boxed{\;Y_i-\widehat g_{-i}(\X_i)=\frac{Y_i-\widehat g(\X_i)}{1-h_{ii}}\;}
\]
Elevando ao quadrado e tirando a média em $i$, sai a fórmula do enunciado. O ganho está
no lado direito: só entra $\widehat g$, ajustado \emph{uma} vez, com todos os dados. Os
$n$ ajustes do LOOCV viraram um.

\smallskip
\textbf{(e)} Pela simetria, a entrada $(i,i)$ de $\bm H\bm H$ é
$\sum_\ell h_{i\ell}h_{\ell i}=\sum_\ell h_{i\ell}^2$; pela idempotência, ela vale
$h_{ii}$. Daí saem as duas desigualdades:
\[
  h_{ii}=\sum_\ell h_{i\ell}^2\;\ge\;0,
  \qquad
  h_{ii}=h_{ii}^2+\sum_{\ell\neq i}h_{i\ell}^2\;\ge\;h_{ii}^2
  \;\Longrightarrow\;h_{ii}(1-h_{ii})\ge0
  \;\Longrightarrow\;h_{ii}\le1 .
\]

\emph{E a hipótese exclui $h_{ii}=1$.} Suponha $h_{ii}=\X_i^\top(\mathbb{X}^\top\mathbb{X})^{-1}\X_i=1$
e tome $\bm v=(\mathbb{X}^\top\mathbb{X})^{-1}\X_i$, que não é nulo (senão $h_{ii}$ seria
$0$). Sem a linha $i$, a matriz de delineamento $\mathbb{X}_{-i}$ satisfaz
$\mathbb{X}_{-i}^\top\mathbb{X}_{-i}=\mathbb{X}^\top\mathbb{X}-\X_i\X_i^\top$, e então
\[
  \mathbb{X}_{-i}^\top\mathbb{X}_{-i}\,\bm v=\X_i-\X_i\,(\X_i^\top\bm v)=\X_i-\X_i\,h_{ii}=\bm0 :
\]
$\mathbb{X}_{-i}$ perderia o posto, contra a hipótese. (A recíproca também vale, e o
exemplo típico é uma coluna indicadora que só a observação $i$ tem: sem ela, a coluna
zera.)

Com $0\le h_{ii}<1$, o fator $1/(1-h_{ii})$ é pelo menos $1$, e (d) dá
\[
  \abs{Y_i-\widehat g_{-i}(\X_i)}\;\ge\;\abs{Y_i-\widehat g(\X_i)}
  \qquad\text{para toda observação } i .
\]
Tirando a média dos quadrados, $\riscohat_{\text{LOO}}$ nunca fica abaixo do erro
quadrático médio de treino. Isso confirma, de um jeito bem concreto, que o erro de
treino é \textbf{otimista} --- e não só em média: em toda amostra, observação por
observação, e com o tamanho da correção escrito.

O $h_{ii}$ mede a \textbf{alavancagem}: o peso da própria $Y_i$ na predição feita em
$\X_i$. Perto de $1$, essa predição é quase só $Y_i$ --- o ajuste passa colado no
ponto, e o resíduo de treino sai pequeno e engana. O fator $1/(1-h_{ii})$ fica grande
justamente para desfazer o engano: sem a observação $i$, as outras quase não informam
o que acontece em $\X_i$. Para ter uma escala,
$\sum_ih_{ii}=\operatorname{tr}(\bm H)=\operatorname{tr}\big((\mathbb{X}^\top\mathbb{X})^{-1}\mathbb{X}^\top\mathbb{X}\big)$
é o número de colunas de $\mathbb{X}$; a alavancagem \emph{média} é, portanto, o número
de parâmetros dividido por $n$, e "alta" quer dizer muito acima disso.

Um último comentário, de método: repare onde cada hipótese trabalhou. A linearidade do
ajuste nas respostas entrou em (a) e (c); o fato de $\widehat g_{-i}$ ser um ajuste de
mínimos quadrados \emph{na mesma classe de funções}, em (b); o posto cheio, na
unicidade de (b) e em $h_{ii}\neq1$; a simetria e a idempotência de $\bm H$, só em (e).
Um método que não tenha alguma dessas propriedades não herda a fórmula
automaticamente.
\end{solucao}

\end{exercicio}

\begin{exercicio}[a mesma validação cruzada, escrita de dois jeitos]
\label{ex:03-duasformas}
% Aula 03 --- as duas formas da fórmula de k dobras (aula x ISLP).
Quem estuda validação cruzada pelas notas de aula e pelo livro encontra duas fórmulas
diferentes para a mesma quantia, e é natural desconfiar de que uma delas esteja
errada. Não está: elas só diferem na unidade que cada uma pondera igualmente.

As notas de aula definem
\[
  \riscohat_{k\text{-CV}}
   = \frac1n\sum_{j=1}^{k}\sum_{i\in L_j}\big(Y_i-g_{-j}(\X_i)\big)^2 ,
\]
e o livro define $\mathrm{CV}_{(k)}=\frac1k\sum_{j=1}^{k}\mathrm{EQM}_j$, com
$\mathrm{EQM}_j=\frac{1}{\abs{L_j}}\sum_{i\in L_j}(Y_i-g_{-j}(\X_i))^2$.
\begin{itens}
  \item \pts{1,2} Reescreva a fórmula das notas de modo a fazer os
        $\mathrm{EQM}_j$ aparecerem, mostrando que
        \[
          \riscohat_{k\text{-CV}}=\sum_{j=1}^{k}\frac{\abs{L_j}}{n}\,\mathrm{EQM}_j .
        \]
  \item \pts{0,8} Olhando para os pesos que você obteve, diga em que condição as
        duas definições coincidem --- e se coincidem exatamente ou apenas
        aproximadamente. Em seguida, nomeie a unidade que cada fórmula pondera
        igualmente.
  \item \pts{0,7} Com $n=50$ e $k=7$, as dobras não podem ter todas o mesmo
        tamanho; suponha que tenham tamanhos $8,8,7,7,7,7,6$. Confirme que somam $50$
        e explique por que, neste caso, os dois números diferem.
  \item \pts{1,3} Suponha que a dobra de $6$ observações tenha $\mathrm{EQM}=10$ e
        todas as outras tenham $\mathrm{EQM}=1$. Calcule os dois estimadores, diga
        qual deles dá mais peso à dobra atípica, e qual você reportaria. Justifique
        a escolha.
\end{itens}

\begin{solucao}
\textbf{(a)} A chave é reconhecer a soma interna. Pela definição de $\mathrm{EQM}_j$,
basta multiplicar dos dois lados por $\abs{L_j}$:
\[
  \sum_{i\in L_j}\big(Y_i-g_{-j}(\X_i)\big)^2=\abs{L_j}\,\mathrm{EQM}_j .
\]
Substituindo na fórmula das notas,
\[
  \riscohat_{k\text{-CV}}
   =\frac1n\sum_{j=1}^{k}\abs{L_j}\,\mathrm{EQM}_j
   =\sum_{j=1}^{k}\frac{\abs{L_j}}{n}\,\mathrm{EQM}_j .
\]
São os \emph{mesmos} $\mathrm{EQM}_j$ do livro, combinados com pesos
$w_j=\abs{L_j}/n$. E esses pesos somam $1$, porque as dobras particionam as $n$
observações: $\sum_j\abs{L_j}=n$. As duas fórmulas são, portanto, médias ponderadas dos
mesmos $k$ números, e só diferem nos pesos: $\abs{L_j}/n$ numa, $1/k$ na outra.

\smallskip
\textbf{(b)} Os pesos coincidem em todas as dobras se e só se $\abs{L_j}/n=1/k$ para todo
$j$, isto é, se todas as dobras têm tamanho $n/k$ --- o que só é possível quando $k$
divide $n$. Nesse caso as duas fórmulas são a mesma soma dos mesmos números com os
mesmos pesos: coincidem \textbf{exatamente}, e não aproximadamente.

A unidade que cada uma pondera igualmente: as notas dão o mesmo peso a cada
\textbf{observação} --- todo resíduo de validação entra com peso $1/n$ ---; o livro dá o
mesmo peso a cada \textbf{dobra}. Visto do lado das observações, o livro dá a cada
resíduo da dobra $j$ o peso $1/(k\abs{L_j})$.

\smallskip
\textbf{(c)} $8+8+7+7+7+7+6=50$.

Os pesos das notas são $8/50$ (duas vezes), $7/50$ (quatro vezes) e $6/50$; os do
livro, $1/7$ para cada dobra. Como os tamanhos diferem, os pesos diferem, e as duas
médias só coincidem por acaso --- quando os $\mathrm{EQM}_j$ são todos iguais, por
exemplo.

Olhando por observação fica mais claro. No livro, cada resíduo recebe
\[
  \frac{1}{7\cdot8}=\frac1{56},\qquad
  \frac{1}{7\cdot7}=\frac1{49}\qquad\text{ou}\qquad
  \frac{1}{7\cdot6}=\frac1{42},
\]
conforme a dobra em que caiu, contra $1/50$ para todos nas notas. Uma observação da
dobra de $6$ pesa $4/3$ de uma da dobra de $8$, só porque divide a sua dobra com menos
gente.

\smallskip
\textbf{(d)} Pelas notas, a dobra atípica entra com peso $6/50$ e as outras $44$
observações com $44/50$:
\[
  \riscohat_{k\text{-CV}}=\frac{6}{50}\times10+\frac{44}{50}\times1=\frac{104}{50}=\mathbf{2{,}08}.
\]
Pelo livro, a dobra atípica é uma das sete, e só:
\[
  \mathrm{CV}_{(7)}=\frac17\big(10+6\times1\big)=\frac{16}{7}\approx\mathbf{2{,}286}.
\]
O livro dá mais peso à dobra atípica: $1/7\approx0{,}143$ contra $6/50=0{,}12$. Ela é a
menor dobra e, ainda assim, vale um sétimo do total.

\emph{Qual reportar?} Para \emph{estimar o risco}, a das notas tem a seu favor um
argumento limpo. O risco é a esperança da perda numa observação nova, e cada resíduo de
validação $\ell_i$ é uma medida dessa perda. Se essas medidas tivessem a mesma média e a
mesma variância $s^2$, e fossem não correlacionadas --- uma idealização, mas é ela que
orienta ---, toda combinação $\sum_iw_i\ell_i$ com pesos somando $1$ teria a mesma
média, e variância $s^2\sum_iw_i^2$. Pela desigualdade de Cauchy--Schwarz,
\[
  1=\Big(\sum_{i=1}^n w_i\Big)^2\le n\sum_{i=1}^n w_i^2 ,
\]
com igualdade só nos pesos iguais, $w_i=1/n$. Pesar as observações igualmente é o que
menos desperdiça informação; dar mais peso a quem caiu numa dobra pequena é deixar o
acaso do corte decidir.

A do livro também se defende. Ela trata cada dobra como uma repetição do experimento
"treinar e validar", e é a escala natural quando o que se quer é comparar as dobras
entre si: é dos $k$ números $\mathrm{EQM}_j$ que sai o erro-padrão da regra de um
erro-padrão. E é a que o \texttt{scikit-learn} calcula ---
\texttt{cross\_val\_score(\dots).mean()} é média não ponderada ---, então o essencial é
\emph{saber qual das duas o seu código calcula} antes de comparar o seu número com o de
outra pessoa.

E a recomendação das notas encerra a discussão pela raiz: com $n$ pequeno, escolha um
$k$ que divida $n$. Aí as duas coincidem, e a pergunta não se coloca.
\end{solucao}

\end{exercicio}

\begin{exercicio}[de que tamanho é o viés da validação cruzada]
\label{ex:03-viescv}
% Aula 03 --- de que tamanho é o viés da validação cruzada, e para que lado.
A validação cruzada estima o risco de um procedimento --- mas não exatamente o
risco que interessa. Há um descompasso embutido, e a boa notícia é que ele tem
sinal conhecido.
\begin{itens}
  \item \pts{1,0} Cada $g_{-j}$ é ajustado com quantas observações? Usando isso,
        argumente que
        $\E\big[\riscohat_{k\text{-CV}}\big]=\risco_{n-n/k}$, onde $\risco_m$
        denota o risco esperado do procedimento treinado com $m$ observações.

        A conta é mais curta do que parece: tome a esperança dentro da soma dupla
        e pergunte-se, para uma parcela qualquer, o que a observação $i$ representa
        para o ajuste $g_{-j}$ que a está prevendo.
  \item \pts{0,8} O alvo que interessa é $\risco_n$, o risco do procedimento
        treinado com a amostra inteira. Supondo que $\risco_m$ decresça com $m$ ---
        mais dados não pioram ---, conclua que a validação cruzada é
        \textbf{pessimista}.
  \item \pts{0,7} Ordene $k=2$, $k=10$ e $k=n$ quanto ao tamanho desse viés,
        justificando pelo número de observações de treino em cada caso.
  \item \pts{1,0} Se o LOOCV tem o menor viés, por que $k=5$ e $k=10$ são os
        valores usuais? Responda pensando em \emph{para que} a validação cruzada
        costuma ser usada.
\end{itens}

\begin{solucao}
\textbf{(a)} O ajuste $g_{-j}$ usa todas as observações menos as da dobra $L_j$,
isto é, $n-\abs{L_j}$ pontos. Com dobras de tamanho $n/k$, são $n-n/k$.

Tome a esperança na definição:
\[
  \E\big[\riscohat_{k\text{-CV}}\big]
   = \frac1n\sum_{j=1}^{k}\sum_{i\in L_j}
     \E\Big[\big(Y_i-g_{-j}(\X_i)\big)^2\Big].
\]
Agora olhe uma parcela isolada. A observação $i$ pertence a $L_j$, e $L_j$ foi
justamente a dobra retirada para treinar $g_{-j}$ --- então $i$ \textbf{não}
participou daquele ajuste. Para $g_{-j}$, ela é uma observação nova e
independente.

Mas "erro esperado de um ajuste numa observação nova" é exatamente a definição de
risco. Como $g_{-j}$ foi treinado com $n-n/k$ pontos, essa esperança vale
$\risco_{n-n/k}$ --- e vale o mesmo para toda parcela, por simetria, já que as
observações são i.i.d.

São $n$ parcelas no total (cada observação aparece em exatamente uma dobra), o
fator $1/n$ cancela, e sobra $\E[\riscohat_{k\text{-CV}}]=\risco_{n-n/k}$.

\smallskip
\textbf{(b)} Como $n-n/k<n$ e $\risco_m$ é decrescente em $m$, temos
$\risco_{n-n/k}\ge\risco_n$. Logo
\[
  \E\big[\riscohat_{k\text{-CV}}\big] = \risco_{n-n/k} \;\ge\; \risco_n .
\]
A validação cruzada superestima, em média, o risco do procedimento treinado com
tudo.

A origem do viés cabe numa frase: \textbf{a CV mede o desempenho de um modelo
treinado com menos dados do que aquele que você vai de fato usar.}

E o sinal conhecido é uma vantagem, não um defeito. Se a CV diz que o modelo é
bom, ele é pelo menos aquilo --- o erro está do lado seguro.

\smallskip
\textbf{(c)} O viés é tanto maior quanto menos observações sobram para treinar:

\begin{center}\small
\begin{tabular}{@{}lll@{}}
\toprule
 & observações de treino & viés \\
\midrule
$k=2$   & $n/2$    & o maior \\
$k=10$  & $0{,}9n$ & pequeno \\
$k=n$ (LOOCV) & $n-1$ & o menor \\
\bottomrule
\end{tabular}
\end{center}

Com $k=2$, metade dos dados é descartada em cada rodada, e o que se estima é o
risco de um procedimento com metade da amostra --- que pode ser bem pior. Com o
LOOCV, cada ajuste usa $n-1$ pontos, praticamente a amostra inteira, e o alvo
estimado é praticamente o certo.

\smallskip
\textbf{(d)} Porque o viés atrapalha menos do que parece para o uso mais comum, e
o LOOCV custa caro.

Se o objetivo for \textbf{reportar} o risco, o viés importa de verdade: o número
sai sistematicamente alto, e aí um $k$ grande ajuda.

Só que o uso mais frequente é \textbf{escolher} um hiperparâmetro --- descobrir
onde a curva $\theta\mapsto\riscohat(\theta)$ tem mínimo. Para isso o que importa
é a \emph{forma} da curva, não o nível dela. Um viés que desloca todos os valores
para cima de maneira parecida preserva a posição do mínimo, e a escolha não muda.

Some a isso o custo: o LOOCV exige $n$ ajustes e, fora do caso linear, não tem
atalho --- justamente o que o exercício do atalho mostra é que a economia só
existe quando o ajuste é linear. Com $k=10$ o viés já é pequeno e o custo é dez
ajustes, não mil.

Uma ressalva, para fechar: o argumento clássico diz também que $k$
grande aumenta a \emph{variância} da estimativa, porque as dobras ficam muito
parecidas entre si. É o que o deck da Aula~03 ensina, e é o que se costuma repetir
--- ainda que a medição desse efeito seja menos consensual do que o argumento
sugere.
\end{solucao}

\end{exercicio}

%==============================================================================
\section*{Aula 04 --- Métodos Não Paramétricos (KNN)}
%==============================================================================

\begin{exercicio}[quanto custa cada nó de um spline]
\label{ex:04-splines}
% Aula 04 --- a base truncada de splines e o custo de cada nó.
Um \emph{spline} é uma função polinomial por partes, colada suavemente nos
\textbf{nós}. A Aula~04 apresenta a base truncada que os constrói, e este
exercício pergunta por que ela funciona --- e quanto custa cada nó.

Um \emph{spline} de grau $g$ com nós $t_1<\dots<t_J$ é combinação linear de
\[
  1,\ x,\ x^2,\ \dots,\ x^{g},\
  (x-t_1)_+^{g},\ \dots,\ (x-t_J)_+^{g},
  \qquad (u)_+=\max(u,0).
\]

\emph{Sobre as letras:} aqui $g$ é o \textbf{grau} e $J$ o \textbf{número de
nós}. As notas usam $k$ e $p$ para essas duas quantidades; trocamos as letras de
propósito, para não colidir com o $k$ do KNN nem com o $p$ de covariáveis. É o
tipo de colisão que a Aula~04 tem e que convém não carregar para uma prova.
\begin{itens}
  \item \pts{0,6} Quantas funções há na base? Escreva $I$ em função de $g$ e $J$,
        e confira no caso cúbico ($g=3$) com $4$ nós.
  \item \pts{1,0} Mostre que $(x-t_j)_+^{g}$ e suas $g-1$ primeiras derivadas se
        anulam em $x=t_j$. Em seguida, tire a conclusão que interessa: por que
        acrescentar essa função à base \emph{não altera} o ajuste à esquerda de
        $t_j$, e por que a emenda no nó sai suave sem que se precise impor
        restrição nenhuma.
  \item \pts{1,0} Use (b) para justificar a frase "cada nó compra exatamente um
        grau de liberdade". Depois compare com a alternativa ingênua: ajustar um
        polinômio de grau $g$ \textbf{independente} em cada um dos $J+1$ pedaços.
        Quantos parâmetros seriam? Onde foi parar a diferença, e o que ela custa?
  \item \pts{0,9} O ajuste é por mínimos quadrados sobre essa base. Isso faz do
        \emph{spline} um \textbf{ajuste linear} no sentido da Aula~02 --- explique
        por quê. E então responda: se ele é linear, o que exatamente o torna um
        método \emph{não paramétrico}?
\end{itens}

\begin{solucao}
\textbf{(a)} São $g+1$ monômios ($1,x,\dots,x^g$) mais $J$ funções truncadas, uma
por nó:
\[
  I = (g+1) + J .
\]
No caso cúbico com $4$ nós, $I=(3+1)+4=\mathbf{8}$ funções na base --- oito
coeficientes a estimar.

\smallskip
\textbf{(b)} Separe os dois lados do nó.

\emph{À esquerda}, $x<t_j$ dá $(x-t_j)_+=0$: a função é identicamente nula num
intervalo inteiro, e uma função identicamente nula tem todas as derivadas nulas
ali.

\emph{À direita}, $x>t_j$ dá $(x-t_j)_+^g=(x-t_j)^g$, cuja $m$-ésima derivada é
\[
  \frac{g!}{(g-m)!}\,(x-t_j)^{g-m}.
\]
Para $m\le g-1$ o expoente $g-m$ é positivo, e essa expressão tende a $0$ quando
$x\downarrow t_j$.

Os limites laterais coincidem, e valem zero, para a função e suas primeiras $g-1$
derivadas. A $g$-ésima é a primeira a saltar: passa de $0$ para $g!$.

Daí as duas conclusões. Como a função é nula à esquerda, somá-la ao modelo
\textbf{não muda nada antes do nó} --- o coeficiente dela é livre para ajustar o
que vem depois sem estragar o que veio antes. E como ela e suas $g-1$ derivadas se
anulam no nó, a emenda é de classe $C^{g-1}$ automaticamente: a suavidade não
precisa ser imposta, ela está embutida na escolha da base.

\smallskip
\textbf{(c)} Cada nó acrescenta \emph{uma} função à base, logo \emph{um}
coeficiente livre, logo um grau de liberdade. E, por (b), esse grau de liberdade é
\textbf{localizado}: só atua à direita de $t_j$. Você compra flexibilidade onde
quer, e paga por uma unidade de cada vez.

A alternativa ingênua seriam $J+1$ pedaços com $g+1$ coeficientes cada:
\[
  (J+1)(g+1)\ \text{parâmetros.}
\]
No caso cúbico com $4$ nós, $5\times 4=20$ --- contra os $8$ da base truncada.

A diferença são as restrições de continuidade. A base truncada já embute que a
função e suas $g-1$ primeiras derivadas são contínuas em cada nó: $g$ restrições
por nó, $gJ$ no total. E a conta fecha:
\[
  (J+1)(g+1)-gJ = g+1+J = I .
\]
Os $12$ parâmetros que sumiram no caso cúbico são exatamente as $3\times 4$
restrições.

O que a alternativa custa: com mais liberdade, ela ajusta \emph{melhor} o treino
--- e produz saltos nos nós, variância alta e extrapolação ruim. Paga-se
flexibilidade com instabilidade, que é a Aula~01 aparecendo de novo com outra
roupa.

\smallskip
\textbf{(d)} Fixadas as funções da base, o modelo é
$g(x)=\sum_{m=1}^{I}\beta_m f_m(x)$: \textbf{linear nos parâmetros} $\beta_m$,
ainda que profundamente não linear em $x$. Basta montar a matriz de delineamento
com $\mathbb{X}_{im}=f_m(x_i)$ e os $\beta_m$ saem pelas equações normais, como na
Aula~02. Vale inclusive a matriz de projeção $\bm H$ --- e, com ela, o atalho do
LOOCV.

O que torna o método não paramétrico não é a forma do ajuste, então. É o fato de
$I$ --- o número de funções da base, isto é, o número de nós --- \textbf{crescer
com $n$}. Com poucos dados usam-se poucos nós e a curva é rígida; com muitos
dados, mais nós, e a classe de funções representáveis se aproxima de qualquer
função suave.

Essa é a distinção que a Aula~04 faz logo no começo: um método paramétrico fixa a
complexidade \emph{antes} de ver os dados; um não paramétrico deixa que ela cresça
com eles. Não é "ter ou não ter parâmetros" --- é \emph{quando} o número deles
fica decidido.
\end{solucao}

\end{exercicio}

\begin{exercicio}[duas doenças, dois remédios]
\label{ex:04-ponte}
% Aula 04 --- a ponte 02<->04: regularizar ataca variância, não paramétrico ataca viés.
As Aulas~02 e~04 apresentam ferramentas que, à primeira vista, andam em sentidos
opostos. A regularização \emph{restringe} o modelo; os métodos não paramétricos o
\emph{soltam}. Parece contradição, e não é: são remédios para doenças diferentes.

Este exercício pede que você organize as duas numa única imagem --- a da Aula~01.
\begin{itens}
  \item \pts{0,7} Escreva a decomposição do erro de um estimador $\rhat$ em
        relação a $r$ nos dois termos da Aula~01, e diga qual dos dois cada
        família de métodos ataca.
  \item \pts{0,8} Descreva, em uma frase cada, a situação que motiva cada família:
        o que está errado com o modelo \emph{antes} de aplicarmos Ridge, e o que
        está errado \emph{antes} de aplicarmos um spline?
  \item \pts{1,0} O KNN com $k=1$ e o MQO com $p$ grande erram por motivos que,
        na decomposição, caem no mesmo lugar. Explique. Em seguida diga qual é o
        remédio de cada um e por que os dois remédios se parecem.
  \item \pts{1,0} As duas ideias podem ser combinadas? Dê um exemplo concreto,
        visto no curso, de método que seja ao mesmo tempo não paramétrico e
        regularizado, e diga qual botão controla cada coisa nele.
\end{itens}

\begin{solucao}
\textbf{(a)} Para um ponto $\x$ fixo,
\[
  \E_\Dados\big[(\rhat(\x)-r(\x))^2\big]
   = \underbrace{\big(\E_\Dados[\rhat(\x)]-r(\x)\big)^2}_{\text{viés}^2}
   + \underbrace{\Var_\Dados\big(\rhat(\x)\big)}_{\text{variância}} .
\]

A \textbf{regularização ataca a variância}: ela restringe o espaço de soluções,
encolhendo os coeficientes, e aceita pagar viés por isso.

Os \textbf{métodos não paramétricos atacam o viés}: eles ampliam a classe de
funções representáveis para que $r$ caiba nela, e aceitam pagar variância.

\smallskip
\textbf{(b)} Antes do Ridge, o problema é um modelo \textbf{flexível demais para a
amostra que se tem}: colunas demais ou colineares, o ajuste persegue o ruído e os
coeficientes ficam instáveis. Excesso de variância.

Antes do spline, o problema é um modelo \textbf{rígido demais para a verdade}: a
reta não acompanha uma $r$ curva, e o erro persiste por mais dados que se colete.
Excesso de viés.

A diferença entre "para a amostra" e "para a verdade" é o coração da distinção:
uma é falta de dados, a outra é falta de classe.

\smallskip
\textbf{(c)} Os dois erram por \textbf{variância}.

O KNN com $k=1$ faz cada predição depender de uma única observação, e todo o ruído
daquele ponto entra na resposta. O ajuste passa exatamente pelos dados de treino,
com resíduo nulo e nenhuma estabilidade --- troque a amostra e a função muda
inteira.

O MQO com $p$ grande tem parâmetros demais para os dados disponíveis, e cada
coeficiente é estimado com pouca informação. No limite $p>n$ a solução nem é
única.

O remédio do KNN é \textbf{aumentar $k$}, fazendo cada predição ser a média de
mais observações. O do MQO é \textbf{regularizar}, encolhendo os coeficientes em
direção a zero. Os dois se parecem porque fazem a mesma coisa no fundo: reduzem a
\emph{complexidade efetiva} do ajuste, trocando variância por viés. Em ambos há um
botão contínuo --- $k$ de um lado, $\lambda$ do outro --- e em ambos ele se
escolhe por validação cruzada, nunca no olho.

Um cuidado que costuma confundir: os botões correm em sentidos opostos. $k$
\emph{grande} e $\lambda$ \emph{grande} dão modelos mais simples, mas $k$ grande
significa vizinhança larga e $\lambda$ grande significa coeficientes pequenos. O
teste seguro, quando a memória falha, é perguntar de que lado está o modelo
constante.

\smallskip
\textbf{(d)} Podem, e o exemplo do curso são as \emph{smoothing splines} da
Aula~04.

Elas põem um nó em \textbf{cada observação} --- portanto são não paramétricas ao
extremo, com $I=n$ --- e controlam a flexibilidade resultante \textbf{penalizando
a curvatura}:
\[
  \sum_{i=1}^n\big(y_i-g(x_i)\big)^2 + \gamma\int \big(g''(x)\big)^2\,dx .
\]

Os dois papéis estão separados e visíveis: a base cresce com $n$ (o lado não
paramétrico) e o $\gamma$ regula quanta daquela liberdade é de fato usada (o lado
regularizado). E neste método quem manda é o $\gamma$ --- os nós deixaram
de ser o botão, porque são todos. É uma solução elegante para um problema chato,
que é ter de escolher onde pôr os nós.
\end{solucao}

\end{exercicio}

%==============================================================================
\section*{Aula 05 --- Árvores de Regressão e \emph{Ensembles}}
%==============================================================================

\begin{exercicio}[podar é outro problema, não o mesmo de crescer]
\label{ex:05-poda}
\input{exercicios/ex-11-poda}
\end{exercicio}

\begin{exercicio}[as duas condições, e o que a floresta faz com elas]
\label{ex:05-condicoes}
\input{exercicios/ex-12-condicoes}
\end{exercicio}

%==============================================================================
\section*{Aula 06 --- Pré-processamento e \emph{Pipelines}}
%==============================================================================

\begin{exercicio}[para quem a régua importa]
\label{ex:06-escala}
% Aula 06 --- quem é sensível à escala das covariáveis, e por quê.
Um colega recebeu um banco em que a coluna \texttt{renda} está em reais e a coluna
\texttt{altura} em metros, e pergunta se precisa padronizar antes de ajustar. A
resposta honesta é "depende do método" --- e o critério que decide cabe numa
pergunta: \textbf{o método soma coisas vindas de colunas diferentes?}

Se soma, as unidades precisam ser comparáveis, e trocar uma delas estraga a
comparação. Se não soma, a unidade é irrelevante.
\begin{itens}
  \item \pts{2,4 --- 0,4 cada} Para cada método abaixo, diga se multiplicar uma
        única coluna por $1000$ \textbf{muda} ou \textbf{não muda} a predição,
        justificando em uma linha onde está --- ou onde não está --- a soma sobre
        colunas.
        \begin{enumerate}[label=(\roman*),leftmargin=1.0cm,nosep]
          \item Mínimos quadrados ordinários.
          \item Regressão Ridge com $\lambda>0$ fixo.
          \item Árvore de regressão.
          \item Floresta aleatória.
          \item KNN com $k$ fixo.
          \item \emph{Spline} de base truncada, ajustado por mínimos quadrados.
        \end{enumerate}
  \item \pts{0,7} Entre os que mudam, qual você espera que seja o mais afetado, e
        por quê? Compare o \emph{modo} como a escala entra em cada um.
  \item \pts{0,9} Um colega observa que, ao reescalar uma coluna, o risco do Lasso
        \emph{melhorou}, e conclui que não padronizar foi uma boa ideia. Explique
        o erro do raciocínio.
\end{itens}

\begin{solucao}
\textbf{(a)}

\begin{center}\small
\renewcommand{\arraystretch}{1.25}
\begin{tabular}{@{}clp{8.2cm}@{}}
\toprule
 & muda? & onde está (ou não) a soma sobre colunas \\
\midrule
(i) & \textbf{não} & O MQO é equivariante: o coeficiente da coluna se divide por
$1000$ e o produto $\beta_jx_j$ fica igual. Nada na função-objetivo compara
coeficientes entre si. \\
(ii) & \textbf{sim} & A penalidade $\lambda\sum_j\beta_j^2$ \emph{soma}
coeficientes de colunas diferentes; mudar a unidade de uma delas muda quanto ela
pesa nessa soma. \\
(iii) & \textbf{não} & A árvore só pergunta "$x_j\le t$?". O corte $t$ é
multiplicado junto, e a pergunta separa exatamente as mesmas observações. \\
(iv) & \textbf{não} & É uma média de árvores, e cada uma é indiferente; a
indiferença atravessa a agregação. \\
(v) & \textbf{sim} & A distância euclidiana $\sum_j(x_{ij}-x_{\ell j})^2$ é uma
soma sobre colunas. \\
(vi) & \textbf{não} & É um MQO sobre uma base fixa de funções de \emph{uma}
coluna; vale o mesmo argumento de (i). \\
\bottomrule
\end{tabular}
\end{center}

\smallskip
\textbf{(b)} O \textbf{KNN}, e a diferença é de ordem, não de grau.

Nos métodos penalizados a escala entra \emph{uma vez}, de forma linear, na
comparação entre coeficientes. Na distância euclidiana ela entra \textbf{ao
quadrado}: multiplicar a coluna por $1000$ multiplica as diferenças naquela
coordenada por $1000$, e a contribuição delas para a distância por $10^6$.

O efeito é que a coluna reescalada passa a decidir sozinha quem é vizinho de quem.
As demais não são apenas menos importantes --- elas somem da conta, e o método
deixa de usar a informação que carregavam.

\smallskip
\textbf{(c)} O erro é confundir um resultado com um argumento.

A melhora aconteceu porque a coluna reescalada era, naquele problema, a de maior
efeito verdadeiro. Deixá-la $1000$ vezes maior tornou o coeficiente dela $1000$
vezes menor, e um coeficiente minúsculo quase não pesa em $\sum_j\abs{\beta_j}$ ---
a coluna ficou praticamente isenta da penalidade. Isentar de penalidade justamente
a coluna que mais importa calhou de ajudar.

Calhou: é sorte, não método. Aplique a mesma reescala a uma coluna
\emph{irrelevante} e o efeito se inverte --- ela ganharia liberdade para receber um
coeficiente grande sem pagar por isso, e o ajuste pioraria.

O ponto de fundo é que, nos dois casos, quem decidiu quanto regularizar cada
coluna não foi o analista: foi a \textbf{unidade de medida} em que o dado chegou.
E é por isso que o argumento do colega não se salva nem quando o resultado o
favorece --- um procedimento cujo desempenho depende de a coluna ter sido anotada
em metros ou em milímetros está fora de controle, e ter dado certo desta vez não
muda isso.
\end{solucao}

\end{exercicio}

\end{document}
